% ======================================================================
%  sec17.tex —— 第 17 节 局部自由模
%  底本：原书 pp.83–84（印刷页 = PDF 页）
%  原书标题：17. Locally free modules
% ======================================================================

\section{局部自由模}\label{sec:17}

\subsection{局部自由模}\label{sec:17.1}

\begin{definition}\label{def:17.1.1}
设 $(X,\cO_X)$ 是环化空间。若对任意 $x\in X$，存在 $x$ 的开邻域 $U\subseteq X$ 与集合 $I$，使 $\cF|_U\simeq\cO_U^{\oplus I}$，则称 $\cO_X$-模 $\cF$ 为\textbf{局部自由}的。

若在上述定义中可以取 $I$ 为有限集，则称局部自由 $\cO_X$-模 $\cF$ 为\textbf{有限秩}的。

若在上述定义中可以取 $I$ 使 $|I|=n$，则称局部自由 $\cO_X$-模 $\cF$ 的\textbf{秩}为 $n$。

秩为 $1$ 的局部自由 $\cO_X$-模也称为\textbf{线丛}（line bundle）。
\end{definition}

\begin{yizhu}
秩可以逐点读出：$\cF$ 局部自由时，函数 $x\mapsto\mathrm{rank}_{\cO_{X,x}}\cF_x$ 在每点邻域内为常数，故在每个连通分支上为常数。原文的「秩为 $n$」要求同一个 $n$ 在每点都能取到，比「逐点秩为 $n$」更强。非连通概形上两者不同：在两个点上分别取秩 $1$ 与秩 $2$ 的自由层，逐点有秩而整体无秩。
[依据：原书 p.83 定义 \ref{def:17.1.1} 的三段表述。]
\end{yizhu}

\begin{yizhu}
「局部自由」一词的来源要在相对谱里看。原书 §\ref{sec:14.3} 把拟凝聚模 $\cE$ 送到 $\mathbf V_S(\cE)=\underline{\Spec}_S(\mathrm{Sym}\,\cE)$；当 $\cE$ 秩 $n$ 局部自由时，$\mathbf V(\cE)$ 局部同构于 $\A^n_S$，正是秩 $n$ 向量丛。Stacks Project 把这一对应表述为向量丛范畴与拟凝聚模范畴之间的反变等价，「秩 $1$ 局部自由」与「线丛」两种叫法由此统一。
[来源：Stacks Project, Tag 01M1（Section 27.6, Vector bundles），\url{https://stacks.math.columbia.edu/tag/01M1}（2026-11-17 访问）。]
\end{yizhu}

下面这个结果是显然的。

\begin{lemma}\label{lem:17.1.2}
设 $X$ 是概形。局部自由的 $\cO_X$-模是拟凝聚的。
\end{lemma}

\begin{lemma}\label{lem:17.1.3}
设 $X$ 是环化空间，$\cL$ 是有限表现的 $\cO_X$-模。设 $x\in X$ 使 $\cL_x\simeq\cO_{X,x}^{\oplus r}$（$r\geq 0$）。则存在 $x$ 的开邻域 $U\subseteq X$，使 $\cL|_U\simeq\cO_U^{\oplus r}$。
\end{lemma}

\begin{proof}
取 $\cL_x$ 作为自由 $\cO_{X,x}$-模的一组基 $s_{1,x},\dots,s_{r,x}$。取 $x$ 的开邻域 $U$，使每个 $s_{i,x}$ 都是截面 $s_i\in\cL(U)$ 的芽。考虑所得态射
\[
\phi\colon\cO_U^{\oplus r}\xrightarrow{(s_1,\dots,s_r)}\cL|_U.
\]
按构造，$\phi_x$ 是同构，故 $\ker(\phi)_x$ 与 $\coker(\phi)_x$ 都为零。

由引理 \ref{lem:16.1.4}(1)，$\coker(\phi)$ 是有限型的；于是由引理 \ref{lem:16.1.6}，其支集 $\mathrm{supp}(\coker(\phi))$ 是 $U$ 的闭子集。因 $x\notin\mathrm{supp}(\coker(\phi))$，可缩小 $U$ 并设 $\coker(\phi)=0$。

另一方面，由习题 \ref{exc:16.2.5}(3)，$\ker(\phi)$ 是有限型的。同法可再缩小 $U$ 并设 $\ker(\phi)=0$。

这就证明了：存在 $x$ 的开邻域 $U$，使 $\phi$ 是同构。
\end{proof}

\begin{yizhu}
这条引理是「一点处同构可传播到邻域」的标准机制，两个条件各用一次。余核由引理 \ref{lem:16.1.4}(1) 知有限型，其支集由引理 \ref{lem:16.1.6} 知闭，故能缩到余核为零；核由习题 \ref{exc:16.2.5}(3) 知有限型，同法缩到核为零。原文把后者误引为「引理 16.2.5(3)」，译文按习题处理。若去掉有限表现，核与余核未必有限型，逐步缩小的论证随即失效。
[依据：原书 p.83 引理 \ref{lem:17.1.3} 的证明。]
\end{yizhu}

\begin{proposition}\label{prop:17.1.4}
设 $X$ 是局部环化空间，$\cL\in\omod$。下述条件等价：
\begin{enumerate}[label=(\roman*),leftmargin=2em,itemsep=0.25ex]
\item $\cO_X$-模 $\cL$ 是秩 $1$ 局部自由的。
\item $\cO_X$-模 $\cL$ 可逆，即存在 $\cO_X$-模 $\cL'$ 使 $\cL\otimes_{\cO_X}\cL'\simeq\cO_X$。
\end{enumerate}
此外，若这些条件成立，则 (ii) 中的 $\cL'$ 同构于 $\cHom_{\cO_X}(\cL,\cO_X)$。
\end{proposition}

\begin{proof}
(i)$\Rightarrow$(ii)：只需说明典范态射
\[
\cL\otimes_{\cO_X}\cHom_{\cO_X}(\cL,\cO_X)\to\cO_X
\]
是同构。该断言在 $X$ 上局部，故可设 $\cL$ 是秩 $1$ 自由模。此时上述态射显然是同构。

(ii)$\Rightarrow$(i)：设 $\phi\colon\cL\otimes_{\cO_X}\cL'\to\cO_X$ 与 $\varphi\colon\cO_X\to\cL'\otimes_{\cO_X}\cL$ 是互逆的同构。注意 $\varphi$ 对应 $\cL'\otimes_{\cO_X}\cL$ 的一个整体截面。由于 (i) 是 $X$ 上的局部性质，可设该截面落在
\[
\cL'(X)\otimes_{\cO_X(X)}\cL(X)\to(\cL'\otimes_{\cO_X}\cL)(X)
\]
的像中。换言之，存在有限多个截面 $s_i\in\cL(X)$ 与 $s'_i\in\cL'(X)$，使
\begin{equation}
\varphi=\sum_{i=1}^{n}s'_i\otimes s_i, \tag{17.1}
\end{equation}
其中把 $s_i$（相应地 $s'_i$）看作态射 $\cO_X\to\cL$（相应地 $\cO_X\to\cL'$）。

考虑同构
\[
\cL\xrightarrow{\id_{\cL}\otimes\varphi}\cL\otimes_{\cO_X}\cL'\otimes_{\cO_X}\cL\xrightarrow{\phi\otimes\id_{\cL}}\cL.
\]
由 (17.1)，上述复合等于
\begin{equation}
\cL\xrightarrow{(\phi(-,s'_1),\dots,\phi(-,s'_n))}\cO_X^{\oplus n}\xrightarrow{(s_1,\dots,s_n)}\cL. \tag{17.2}
\end{equation}
由此得 $\cO_X^{\oplus n}\simeq\cL\oplus\cM$，其中 $\cM$ 是 (17.2) 中第二个映射的核。于是有正合列
\[
\cO_X^{\oplus n}\to\cO_X^{\oplus n}\to\cL\to 0,
\]
其中第一个映射是复合 $\cO_X^{\oplus n}\to\cM\to\cO_X^{\oplus n}$。特别地，$\cL$ 是有限表现的。

现在对任意点 $x\in X$ 有 $\cL_x\otimes_{\cO_{X,x}}\cL'_x\simeq\cO_{X,x}$，这说明 $\cL_x$ 是局部环 $\cO_{X,x}$ 上的可逆模。交换代数中的熟知事实指出，此时 $\cL_x$ 是秩 $1$ 自由 $\cO_{X,x}$-模。由引理 \ref{lem:17.1.3}，存在 $x$ 的开邻域 $U\subseteq X$，使 $\cL|_U$ 是秩 $1$ 自由 $\cO_U$-模。这就证明了 (i)。

剩下要证：当 (i) 与 (ii) 成立时 $\cL'\simeq\cHom_{\cO_X}(\cL,\cO_X)$。这由命题 \ref{prop:12.4.10} 得到。
\end{proof}

\begin{yizhu}
命题假设 $X$ 是局部环化空间，这一条不能省：只有茎是局部环时，$\cL_x\otimes\cL'_x\simeq\cO_{X,x}$ 才能推出 $\cL_x$ 秩 $1$ 自由。一般环化空间上逆命题不成立，反例可以取得很小：取单点空间 $X$ 与 $\cO_X=\Z[\sqrt{-5}]$，非主理想 $(2,1+\sqrt{-5})$ 可逆却非自由，故给出可逆但不局部自由的层。
[来源：Stacks Project, Tag 0B8M（Lemma 17.25.4，秩 $1$ 局部自由与可逆的关系），\url{https://stacks.math.columbia.edu/tag/0B8M}（2026-11-17 访问）。]
\end{yizhu}

\begin{yizhu}
证明中「可设该截面落在像中」一句是局部化的简写，严格说须先取开覆盖。$\varphi$ 是 $\cL'\otimes\cL$ 的整体截面，按张量积层的构造，它在每个足够小的开集上可写成有限和 $\sum s'_i\otimes s_i$；而 (i) 是局部性质，故可在这样的开集上论证。原文略去了这一步，读者宜补上。
[译者按（上下文推断，无外部依据）：原文 p.84 第 5–10 行把「层化」的取覆盖步骤压缩成了一句话。]
\end{yizhu}

\begin{yizhu}
末尾的 $\cL'$ 就是 $\cL$ 的对偶：由 $\cL\otimes\cL'\simeq\cO_X$ 与张量积、内 Hom 的伴随关系得 $\cL'\simeq\cHom_{\cO_X}(\cL,\cO_X)$。故秩 $1$ 局部自由模的对偶仍秩 $1$ 局部自由，且 $\cHom(\cL,\cO_X)\otimes\cL\simeq\cO_X$；这正是 Picard 群的群运算，$(ii)$ 于是把线丛的同构类组织成一个群。反过来，Picard 群的每个元素都由一个线丛的同构类代表，两种说法彼此对应。
[依据：原书 p.84 命题 \ref{prop:17.1.4} 末段与其所引命题 \ref{prop:12.4.10}。]
\end{yizhu}

